\section{Structural Fixing and the Higgs Interpretation}
\label{sec:higgs}

\subsection{Reformulating the Higgs Mechanism}

In the Standard Model, the Higgs mechanism is introduced to explain
symmetry breaking and mass generation.
The Higgs field is treated as an additional entity imported for this purpose.
Within the present framework, the question is reformulated:
\begin{equation}
  \boxed{
    \text{Can symmetry breaking and mass arise from internal
    structural dynamics?}
  }
\end{equation}

\subsection{Fluctuating Degenerate Structure}

Section~\ref{sec:weak} showed that degenerate triangular closure
admits internal fluctuation between two states:
\begin{equation}
  |L\rangle \;\leftrightarrow\; |R\rangle.
\end{equation}
In this regime, no configuration is uniquely preferred
and the structure remains dynamically unstable.
\begin{equation}
  \boxed{\text{The degenerate closure is intrinsically undecided.}}
\end{equation}

\subsection{Order Parameter for Structural Asymmetry}

To characterize this instability, introduce an order parameter:
\begin{equation}
  \label{eq:higgs_order}
  \varphi = |\psi_L|^2 - |\psi_R|^2.
\end{equation}
This quantity measures the imbalance between the two internal states.
\begin{equation}
  \boxed{\varphi \;\text{quantifies structural asymmetry in degenerate closure.}}
\end{equation}

\subsection{Fixing of Asymmetry}

Stability is achieved when the system selects a preferred configuration,
$\varphi \neq 0$.
Fluctuation is suppressed, a dominant state emerges,
and the structure stabilizes.
\begin{equation}
  \boxed{\text{Fixing of asymmetry leads to stable closure.}}
\end{equation}
This process corresponds to what is conventionally described as
spontaneous symmetry breaking.

\subsection{Relation to Mass}

As asymmetry becomes fixed,
closure becomes more stable and the closure deficit $\Delta$ decreases.
Given equation~\eqref{eq:mass}:
\begin{equation}
  \boxed{
    \text{Mass can be interpreted as a consequence of stabilized asymmetry.}
  }
\end{equation}

\subsection{Effective Potential}

The dynamics of the order parameter are described by an effective potential:
\begin{equation}
  \label{eq:higgs_potential}
  V(\varphi) = a\varphi^2 + b\varphi^4 - h\varphi.
\end{equation}
For $a < 0$, a nonzero minimum emerges:
\begin{equation}
  \langle\varphi\rangle = v_\varphi \neq 0.
\end{equation}
\begin{equation}
  \boxed{\text{A nonzero expectation value corresponds to structural fixing.}}
\end{equation}
The linear term $h$ reflects residual asymmetry intrinsic to the closure structure ---
specifically, the imbalance between left and right edges
and the causal flow asymmetry of the degenerate triangle.

\subsection{Structural Interpretation of the Higgs Field}

Within this framework, the Higgs is not introduced as an external field,
but is associated with an internal degree of freedom
already present in the degenerate closure.
\begin{equation}
  \boxed{
    \text{The Higgs can be interpreted as the order parameter governing
    structural asymmetry.}
  }
\end{equation}

\subsection{Vacuum Expectation Value}

The vacuum expectation value $v_\varphi$ characterizes the stable configuration:
\begin{equation}
  \label{eq:vev}
  v_\varphi \sim \sqrt{
    \frac{\alpha_1 U(A) + \alpha_2\,\delta_{\mathrm{asym}}(A) - a_0}
    {2\kappa_B}
  }.
\end{equation}
This expression encodes non-closure contributions, asymmetry contributions,
and structural constraints.
\begin{equation}
  \boxed{
    \text{The vacuum state corresponds to a stabilized closure configuration.}
  }
\end{equation}

\subsection{Structural Inversion}

The conventional interpretation is:
$\text{symmetry breaking} \;\rightarrow\; \text{mass}.$
Within this framework:
\begin{equation}
  \boxed{\text{structural fixing} \;\rightarrow\; \text{mass.}}
\end{equation}
Mass is not introduced through an external mechanism,
but emerges from stabilization of internal structure.

\subsection{Summary of Vol.\ 2}

The structural correspondence developed across this volume:
\begin{equation}
  \boxed{
    \begin{aligned}
      \text{Particle}          &= \text{triangular closure} \\
      \text{Symmetry}          &= \text{invariance of closure} \\
      \text{Generation}        &= \text{sequential stabilization} \\
      \text{Charge}            &= \text{phase asymmetry} \\
      \text{Weak interaction}  &= \text{structural fluctuation} \\
      \text{Mass}              &= \text{stabilized asymmetry}
    \end{aligned}
  }
\end{equation}
